\chapter{Introduction}
Virtual reality (VR) enables the creation of immersive, interactive environments that elicit a strong sense of presence by enveloping users in computer-generated worlds \cite{slater2018immersive,freeman2017virtual}. This unique capability has led to its adoption not only in entertainment, but also in education, training, therapy, and healthcare applications \cite{baker2019vrtherapy}. Virtual reality has been increasingly seen as a tool for strengthening mental health and wellness because it provides protected and distraction-free spaces where users can retreat from external stressors and engage in relaxation practices \cite{rehm2021virtual,bailenson2020experience}.

In conjunction with such advances in technology, meditation has traditionally been used as a tool to develop relaxation, improve emotional regulation, and reduce anxiety and stress\cite{sedlmeier2012meta,lutz2008attention,manocha2011mind}. Scientific evidence illustrates the ability of meditation to enhance psychological resilience and promote general mental health. Extensive research shows that meditation can reduce stress and anxiety, improve concentration, and overall well-being \cite{goyal2014meditation,zeidan2015mindfulness}. However, many people struggle to maintain focus during meditation due to distractions, lack of guidance, difficulty staying focused, or inability to relax the mind.

VR offers a very useful bridge; it can turn proven relaxation and mindfulness methods into engaging sensory experiences that make it easier for people to start and keep practicing. The combination of the two spaces, meditation in VR spaces, has created \textbf{VR-centered meditation applications}. These systems present serene and immersive environments that can facilitate the user's ability to transcend peripheral interruptions and focus more intensely on their meditation \cite{liszio2018nature,weibel2020biofeedback}. However, most commercial VR meditation solutions are unchanging, presenting pre-programmed content that is not dynamically adjusted based on the user's state of emotion or physiology. This static responsiveness can inhibit their ability to facilitate increased relaxation and prolonged focus. In order to fill the gap, this thesis is motivated by the opportunity to design and evaluate an \emph{interactive, AI-powered, voice-interactive VR meditation system} which can guide users through meditation based on user stress level. 

\section{Topic Description}
This thesis investigates a \textbf{voice-interactive, AI-guided VR meditation} system developed in Unity for HTC VIVE Pro. The application delivers \emph{natural-language} voice commands (generated by a GPT) to guide breathing, attention control, and relaxation strategies, while \emph{adaptive audiovisual transitions} (e.g. brightening and fading of light during eyes-closed phases to help with breathing, along with meditation music) shape the meditative ambience.

To ground personalisation, the user will be asked about their stress level and the guidance will be given according to that level. The system also records physiological signals \textbf{(SCL) Skin Conductance Level}, \textbf{(SCR) Skin Conductance Response}, \textbf{Skin Temperature}, \textbf{(BVP) Blood Volume Pulse}, \textbf{(PVA) Pulse Volume Amplitude}, \textbf{Pulse / Heart Rate}, \textbf{Motion / Motility} and \textbf{Heart Rate variability (HRV)} \textit{derived from BVP}—using Biofeedback 2000 Xpert sensors. These data are collected during the session and later is used for the comparison of the outcome of meditation \cite{lehrer2013hrv,chittaro2018placebo}. A user study with 44 participants compares the AI-guided VR meditation against a conventional (non-VR) meditation baseline, combining physiological outcomes with self-reports of relaxation and affect.

\section{Research Questions and Hypotheses}
\textbf{RQ1.} \emph{Does AI-guided voice interaction combined with adaptive audiovisual cues improve relaxation and mindfulness during VR meditation compared to a conventional meditation?}  

\textbf{H1.} Participants testing the AI-guided VR meditation will exhibit significantly greater post-session relaxation (physiological and self-report) than conventional meditation. 

\textbf{RQ2.} \emph{Do “eyes-closed mode” visuals (e.g., breathing light effect presented during instructed eyes-closed phases) enhance the perceived depth of meditation and calmness?}  

\textbf{H2.} Participants will report higher immersion and calmness during eyes-closed phases augmented by soothing visuals and ambient audio.

\textbf{RQ3.} \emph{Can a combination of physiological (HRV, skin conductance etc) and self-report measures provide valid, convergent indicators of meditation effectiveness in an AI-guided VR environment?}  

\textbf{H3.} Physiological improvements (e.g., increased HRV, reduced electrodermal activity) will correlate with higher self-reported relaxation and positive affect, indicating convergent validity.


\section{Scientific Contributions}
\begin{itemize}
  \item \textbf{Design Contribution:} A multimodal, \emph{AI-guided voice} and \emph{audiovisual-adaptive} VR meditation prototype that personalizes ambience and guidance to support relaxation, building on best practices in user-centered VR design.
  \item \textbf{Methodological Contribution:} An evaluation protocol combining \emph{physiological} (HRV, skin conductance) and \emph{psychological} (self-report) measures to assess meditative outcomes in VR.
  \item \textbf{Empirical Contribution:} Evidence on whether AI-guided voice prompts and adaptive audiovisual transitions improve relaxation relative to a conventional meditation baseline, including analysis of “eyes-closed mode” visuals.
  \item \textbf{Synthesis Contribution:} A literature-based rationale connecting meditation science, biofeedback, and adaptive interaction in VR, clarifying where adaptive intelligence adds value beyond static experiences.
\end{itemize}

\section{Chapter Outline}
The remainder of this thesis is organised as follows.  
\textbf{Chapter 2} reviews related work on virtual reality meditation, biofeedback-based stress regulation, AI-driven adaptivity, and multimodal interaction design, and identifies the research gap that motivates the present study.  
\textbf{Chapter 3} describes the methodology, including the experimental design, hardware and software architecture, system implementation, participant procedures, measurement instruments, and data processing pipeline.   
\textbf{Chapter 4} presents the empirical results, covering both physiological responses and questionnaire outcomes, and compares the effectiveness of conventional versus AI-guided VR meditation.   
\textbf{Chapter 5} synthesises and interprets the findings in relation to the research questions, discusses their implications, and summarises the overall contributions of the study. 
\textbf{Chapter 6} outlines limitations and proposes future research directions for adaptive, 
biofeedback-aware VR meditation systems.
